%% Allgemeines -----------------------------------

% Fortlaufende Nummerierung - Entweder 1, 2, 3, 4, ... oder mit einem semantischen Schema
\def \decisionID {12}

% Datum, an dem der Record anefertigt wurde
\def \decisionDate {26.11.2026}

% Wer hat an dem Record mitgewirkt?
\def \decisionParticipants {Quentin Becker}

% Der Status in dem sich die Anforderung in diesem Record aktuell befindet.
% Möglichkeiten: Vorgeschlagen, Akzeptiert, Abgelehnt, Abgelöst
\def \decisionStatus {Abgelöst}

% ID des Records, der diesen Record ablöst, falls zutreffend
\def \decisionSuperseded {20}

% ID des Records, den dieser Record ablöst, falls zutreffend
\def \decisionSuperseding {}

%% Anforderung -----------------------------------

% Titel der Anforderung
\def \requirementTitle {Mit VisualBasic auf die XML Struktur der .pptx Dateien zugreifen}

% Kontext der Anforderung
\def \requirementContext {Alle Formate der Microsoft Office Suite basieren auf einem zugrundeliegenden XML-Dokument. In diesem werden die Textabschnitte, Graphiken und Bilder gespeichert. Die Office-Programme interpretieren diesen Inhalt dann zu einem Dokument.}

% Konsequenz der Anforderung
\def \requirementConsequences {Similar Slides kann direkt auf die innere Struktur der Präsentation zugreifen und muss keine rein visuelle Analyse von außen durchführen. Mit der Schnittstelle zwischen unserem Programm und XML ist unsere Einschätzung und Berechnung der Ähnlichkeit präziser, schneller und verlässlicher.}
